\chapter{Gum Infection}

\section*{Definition and Overview}
A gum infection is an infection of the gums and surrounding tissues of the mouth, most commonly caused by the bacteria in dental plaque. Gum infections range from gingivitis, a mild reversible inflammation, to periodontitis, which damages the tissues and bone supporting the teeth, to acute infections such as periodontal abscesses and necrotising ulcerative gingivitis, which cause pain, swelling, and pus. Gum infections can also spread, causing cellulitis of the face, and rarely serious complications. Treatment involves professional dental cleaning, antibiotics for acute infections, improved oral hygiene, and, for abscesses, drainage. Most gum infections are preventable and treatable.

\section*{Epidemiology}
Gum infections are very common. Gingivitis affects most people at some stage, and moderate to severe periodontitis affects around 20--30\% of Australian adults. Acute gum infections, such as periodontal abscesses, are a common reason for emergency dental visits. Necrotising ulcerative gingivitis is less common and mainly affects young adults who smoke, are stressed, or have poor nutrition and oral hygiene. Risk of gum infection is increased by smoking, diabetes, poor oral hygiene, and immune suppression, and it increases with age.

\section*{Aetiology and Pathophysiology}
Gum infections are caused by bacteria that accumulate in plaque and invade the gum tissues.

\begin{itemize}
    \item \textbf{Plaque bacteria:} the mixed bacterial community in dental plaque is the main cause
    \item \textbf{Inflammation:} bacteria trigger an inflammatory response in the gums
    \item \textbf{Gingivitis:} inflammation confined to the gums, causing bleeding and swelling
    \item \textbf{Periodontitis:} infection spreads below the gumline, destroying the attachment and bone
    \item \textbf{Periodontal abscess:} pus collects in a deep pocket, causing a painful, localised swelling
    \item \textbf{Necrotising infection:} tissue-destroying bacteria cause painful ulceration, typically in smokers or stressed individuals
    \item \textbf{Spread:} infection can track into the deeper tissues of the face and neck (cellulitis, Ludwig angina)
    \item \textbf{Risk factors:} smoking, diabetes, stress, poor nutrition, and immune suppression
\end{itemize}

\section*{Clinical Features}
Symptoms depend on the type and severity of the infection.

\begin{itemize}
    \item Bleeding, red, or swollen gums
    \item Persistent bad breath and bad taste
    \item Pain, tenderness, or a throbbing feeling in the gums
    \item A localised swelling or "gum boil" (periodontal abscess)
    \item Pus draining from between the gum and tooth
    \item Loose teeth or a change in the bite
    \item In necrotising gingivitis: painful, bleeding ulcers at the tips of the gum between the teeth, with a foul odour
    \item Fever and feeling unwell with acute infections
    \item Spreading facial swelling in severe cases
\end{itemize}

\section*{Differential Diagnosis}
Other conditions can mimic gum infection.

\begin{itemize}
    \item Dental abscess from an infected tooth root (periapical abscess)
    \item Pericoronitis, infection around an erupting wisdom tooth
    \item Herpes gingivostomatitis and other viral infections
    \item Oral thrush (candidiasis)
    \item Oral cancer, which can present as a non-healing ulcer or swelling
    \item Trauma to the gums
    \item Allergic reactions to toothpaste or other products
    \item Medication-related gum changes
    \item Systemic conditions causing gum inflammation, such as leukaemia
\end{itemize}

\section*{When to Seek Medical Care}
See a dentist promptly for pain, swelling, pus, or bleeding that is persistent or worsening. People with diabetes, immune suppression, or spreading facial swelling should seek care urgently.

\textbf{Call 000 (or local emergency services) for:}
\begin{itemize}
    \item Spreading swelling of the face, neck, or under the jaw with difficulty breathing or swallowing
    \item High fever with facial swelling
    \item Uncontrolled bleeding from the gums
    \item Signs of sepsis, including confusion, rapid breathing, or collapse
\end{itemize}

\section*{Diagnosis and Investigations}
Gum infections are diagnosed by dental examination, with imaging when needed.

\begin{itemize}
    \item \textbf{Clinical examination:} inspection of the gums, probing of pockets, and assessment of loose teeth
    \item \textbf{Dental X-rays:} show bone loss and abscesses around tooth roots
    \item \textbf{Probing:} measures pocket depth
    \item \textbf{Swabs and cultures:} taken from draining pus in some cases
    \item \textbf{Blood tests:} for fever or suspected spread, including inflammatory markers
    \item \textbf{Review of risk factors:} diabetes, smoking, and medicines
\end{itemize}

\section*{Management}
Treatment removes the infection, controls inflammation, and restores oral health.

\begin{itemize}
    \item \textbf{Professional cleaning:} scaling and root planing remove plaque and tartar from above and below the gumline
    \item \textbf{Drainage:} a periodontal abscess is incised and drained
    \item \textbf{Antibiotics:} oral antibiotics for acute or spreading infections; topical antiseptics for local disease
    \item \textbf{Improved oral hygiene:} correct brushing, flossing, and interdental cleaning
    \item \textbf{Antiseptic mouthwash:} chlorhexidine rinses for short-term use
    \item \textbf{Pain relief:} paracetamol or ibuprofen
    \item \textbf{Smoking cessation:} essential for healing and prevention
    \item \textbf{Follow-up:} review of healing and periodontal maintenance visits
    \item \textbf{Surgery:} for persistent deep pockets or advanced periodontitis
\end{itemize}

\section*{Complications}
\begin{itemize}
    \item Spread of infection into the face and neck (cellulitis)
    \item Ludwig angina, a rare but life-threatening infection of the floor of the mouth
    \item Tooth loss from advanced periodontitis
    \item Abscess formation and recurrence
    \item Bone loss around the teeth
    \item Spread of bacteria to the bloodstream, which is a particular risk for people with heart valve conditions
    \item Worsening of diabetes control
    \item Chronic bad breath and effects on quality of life
\end{itemize}

\section*{Prognosis and Outcomes}
With prompt and appropriate treatment, most gum infections resolve fully. Gingivitis is reversible, abscesses settle with drainage and antibiotics, and periodontitis can be controlled with ongoing maintenance. Necrotising gingivitis responds quickly to cleaning, antibiotics, and hygiene improvements. Serious complications are uncommon when treatment is timely. The keys to a good outcome are early dental care, treating underlying risk factors such as smoking and diabetes, and maintaining excellent oral hygiene after the infection clears.

\section*{Prevention and Self-Care}
\begin{itemize}
    \item Brush twice daily and clean between the teeth daily
    \item See a dentist regularly, and promptly for any gum symptoms
    \item Do not smoke
    \item Manage diabetes well
    \item Manage stress and maintain good nutrition
    \item Rinse with warm salt water for comfort after treatment
    \item Complete any prescribed course of antibiotics
    \item Attend follow-up visits and periodontal maintenance appointments
    \item Seek urgent care for spreading swelling, fever, or difficulty swallowing
\end{itemize}

\section*{Sources and Further Reading}
The following reputable sources provide further information for healthcare students and patients seeking reliable, current advice about gum infection.

\begin{itemize}
    \item Australian Dental Association. \textit{Dental abscess and gum infection information.} https://www.ada.org.au/
    \item Healthdirect Australia. \textit{Gum disease and infections.} https://www.healthdirect.gov.au/
    \item Dental Health Services Victoria. \textit{Gum infections and treatment.} https://www.dhsv.org.au/
    \item NSW Health. \textit{Oral health and gum disease.} https://www.health.nsw.gov.au/
    \item NHS. \textit{Gum disease and dental abscesses.} https://www.nhs.uk/
    \item Mayo Clinic. \textit{Periodontal abscess and gum infections.} https://www.mayoclinic.org/
\end{itemize}
